\begin{textsample}{POS Dim 1 – vem\_ed\_22 – Score 26.00 – vem\_ed\_22/t115.txt}  \label{ex:f1_pos_172}
\textbf{continuação} Fatec Itaquaquecetuba - Connecting Traditions and Innovations: The Assistance of Information Technology in the Preservation of Brazilian Indigenous Culture, dos estudantes Lorena Araújo Pereira e Gabriel Nobrega Monteiro, orientados por Wilton Garcia; Fatec Itaquaquecetuba - Earth Timer: proposal for a digital application, do estudante Patrick Paiva, orientado por Wilton Garcia; Fatec Pompeia – Case study on water lackage in the northeast of Brazil, do estudante Israel Nascimento de Oliveira, orientado por Vânia Regina Alves de Souza (leia depoimento da professora a \textbf{seguir}). “Case study on water lackage in the northeast of Brazil trata-se de um estudo de caso realizado em um PCI entre Fatec Pompeia e a Yerevan State University (Armênia), no qual investigamos as principais razões causadoras da falta de água no Sertão no Nordeste brasileiro e buscamos explorar \textbf{possíveis} soluções para enfrentar esse desafio, atendendo aos objetivos de Desenvolvimento \textbf{Sustentável} \textbf{propostos} pela ONU.
O \textbf{intercâmbio} virtual desempenha um \textbf{papel} crucial no \textbf{desenvolvimento} das \textbf{competências} \textbf{linguísticas} e interculturais dos alunos, oferecendo uma plataforma para a prática autêntica da língua estrangeira em \textbf{contextos} reais de comunicação.
Além disso, \textbf{promove} a exposição a diferentes culturas, \textbf{ampliando} a \textbf{compreensão} e o respeito pela diversidade.
Essa experiência virtual \textbf{permite} aos alunos interagir com falantes nativos e outros estudantes ao redor do mundo, enriquecendo seu \textbf{aprendizado} de \textbf{forma} \textbf{significativa} e preparando-os para um mundo globalizado e multicultural.
Participar de um evento em uma universidade estrangeira \textbf{foi} uma experiência enriquecedora \textbf{que} agregou imenso valor tanto para mim, como professora de língua estrangeira, \textbf{quanto} para meus alunos.
A oportunidade de apresentar um trabalho em um \textbf{ambiente} acadêmico internacional proporcionou uma \textbf{troca} de \textbf{conhecimentos} e experiências únicas.
Para mim, \textbf{foi} uma chance de aprimorar minhas \textbf{habilidades} de comunicação e pesquisa, além de \textbf{ampliar} minha rede de contatos \textbf{profissionais}.
Para meus alunos, essa \textbf{vivência} \textbf{foi} inspiradora e motivadora, mostrando-lhes a \textbf{importância} da pesquisa e da participação \textbf{ativa} na comunidade acadêmica global.
Além disso, apresentar um trabalho em um Conclave internacional não só valorizou o esforço e a dedicação dos alunos, mas também os preparou para desafios \textbf{futuros}, \textbf{fortalecendo} suas \textbf{habilidades} de apresentação e confiança no uso da língua estrangeira em \textbf{contextos} acadêmicos e \textbf{profissionais}.”

% matched lemmas: ambiente, ampliar, aprendizado, ativo, competência, compreensão, conhecimento, contexto, continuação, desenvolvimento, forma, fortalecer, futuro, habilidade, importância, intercâmbio, linguístico, papel, permitir, possível, profissional, promover, propor, quanto, que, seguir, ser, significativo, sustentável, troca, vivência
\end{textsample}
